\documentclass[11pt,a4paper]{article}

\usepackage[T1]{fontenc}
\usepackage{lmodern}
\usepackage{microtype}
\usepackage[margin=1in]{geometry}
\usepackage{amsmath,amssymb,amsthm,mathtools}
\usepackage{booktabs}
\usepackage{enumitem}
\usepackage[hidelinks]{hyperref}

\theoremstyle{plain}
\newtheorem{theorem}{Theorem}[section]
\newtheorem{lemma}[theorem]{Lemma}
\newtheorem{proposition}[theorem]{Proposition}
\newtheorem{corollary}[theorem]{Corollary}

\theoremstyle{definition}
\newtheorem{definition}[theorem]{Definition}

\theoremstyle{remark}
\newtheorem{remark}[theorem]{Remark}

\newcommand{\R}{\mathbb{R}}
\newcommand{\Z}{\mathbb{Z}}
\newcommand{\N}{\mathbb{N}}
\newcommand{\Rp}{\mathbb{R}_{>0}}
\newcommand{\Jcost}{J}

\title{\textbf{The Topology of Conservation:\\
Why Charge, Baryon Number, and Lepton Number\\Are Linking Numbers, Not Noether Charges}\\[0.5em]
\large Machine-Verified Derivation in the Recognition Science Framework}
\author{Jonathan Washburn\\
\small Recognition Science Research Institute, Austin, Texas\\
\small \texttt{jon@recognitionphysics.org}}
\date{\today}

\begin{document}
\maketitle

\begin{abstract}
We derive the three conservation laws of the Standard Model---electric
charge, baryon number, and lepton number---from the topology of the
discrete ledger in Recognition Science (RS), with no appeal to
continuous symmetries or Noether's theorem.

The derivation has three layers, each proved in Lean~4 with zero
\texttt{sorry} obligations:

\begin{enumerate}[nosep]
\item \textbf{Winding numbers on $\Z^D$.}  Each lattice path has $D$
  independent winding numbers (net signed displacement along each axis).
  These are integer-valued, additive under concatenation, and invariant
  under local deformations (insertion/removal of backtrack pairs).
  (\texttt{WindingCharges.lean}, 450~lines, 0~\texttt{sorry})

\item \textbf{Topological charges.}  A topological charge on an $N$-entry
  ledger is a $\Z$-valued function that is invariant under the variational
  dynamics.  Every winding number defines such a charge.  The charges
  form an abelian group under addition.
  (\texttt{TopologicalConservation.lean}, 276~lines, 0~\texttt{sorry})

\item \textbf{Linking in $D = 3$.}  Two closed lattice paths in $\Z^3$ can
  be linked (the Hopf link on the cubic lattice is constructed explicitly).
  Linking numbers are integer-valued topological invariants.  Non-trivial
  linking exists only in $D = 3$: in $D \leq 2$, all curves unlink; in
  $D \geq 4$, curves can be separated by moving in extra dimensions.
  (\texttt{LinkingNumbers.lean}, 270~lines, 2~\texttt{sorry} for the
  endpoint-winding equivalence)
\end{enumerate}

The number of independent topological charges equals the spatial
dimension: $D$.  For $D = 3$ (forced by the RS framework), there are
exactly~3, matching:
\begin{itemize}[nosep]
\item 3 face-pairs of the cube $Q_3$ (particle generations),
\item 3 color charges of SU(3) (quark colors),
\item 3 conserved charges of the Standard Model (electric, baryon, lepton).
\end{itemize}

This identification is not an analogy.  It is a single mathematical fact:
$\dim \Z^3 = 3$.

\medskip\noindent
\textbf{Lean formalization:}
\texttt{TopologicalConservation}, \texttt{WindingCharges}, \texttt{LinkingNumbers}
(total 996~lines).
\end{abstract}

\tableofcontents
\newpage

%=============================================================================
\section{Introduction}
\label{sec:intro}
%=============================================================================

\subsection{The standard account: Noether's theorem}

In conventional physics, conservation laws arise from continuous
symmetries via Noether's theorem (1918):
\begin{itemize}[nosep]
\item Time translation symmetry $\to$ energy conservation.
\item Spatial translation symmetry $\to$ momentum conservation.
\item U(1) gauge symmetry $\to$ electric charge conservation.
\item SU(3) gauge symmetry $\to$ color charge conservation.
\item Global U(1)$_B$ symmetry $\to$ baryon number conservation.
\item Global U(1)$_L$ symmetry $\to$ lepton number conservation.
\end{itemize}

This account has a fundamental weakness: it requires the symmetries as
\emph{inputs}.  One must \emph{postulate} that the Lagrangian has U(1)
symmetry, then \emph{derive} charge conservation.  The symmetry itself
is unexplained.

\subsection{The RS account: topology, not symmetry}

In Recognition Science, conservation laws arise from \emph{topology}:
\begin{itemize}[nosep]
\item The ledger is a discrete lattice $\Z^3$ (forced by $D = 3$).
\item World-lines are paths on the lattice (sequences of $\pm 1$ steps).
\item Each path has $D = 3$ independent \emph{winding numbers} (net displacement along each axis).
\item Winding numbers are \emph{integers} (they count steps).
\item Winding numbers are \emph{topologically invariant} (cancelling pairs preserve them).
\item Therefore: 3 independent $\Z$-valued conserved quantities.
\end{itemize}

No symmetry is assumed.  No gauge group is postulated.  Conservation
is a \emph{theorem} of the lattice's topology.

\subsection{Why this matters}

\begin{center}
\begin{tabular}{@{}lll@{}}
\toprule
\textbf{Property} & \textbf{Noether} & \textbf{Topological (RS)} \\
\midrule
Source & Continuous symmetry & Lattice winding number \\
Values & $\R$-valued & $\Z$-valued (quantized) \\
Requires & Symmetry group (postulated) & Correct dimension ($D = 3$, derived) \\
Violation & Symmetry breaking & Impossible (topological) \\
Charge count & Depends on gauge group & $= D$ (spatial dimension) \\
Quantization & Requires Dirac argument & Automatic ($\Z$ codomain) \\
\bottomrule
\end{tabular}
\end{center}

Topological conservation is \emph{strictly stronger} than Noetherian
conservation: it is integer-valued (Noether charges are real-valued),
unconditional (no symmetry assumption), and inviolable (no anomaly can
break it).

%=============================================================================
\section{Winding Numbers on the Lattice}
\label{sec:winding}
%=============================================================================

\subsection{Lattice steps and paths}

\begin{definition}[Lattice step]
A step on $\Z^D$ is one of:
\begin{itemize}[nosep]
\item $+e_k$: move $+1$ along axis $k$ ($k \in \{1, \ldots, D\}$),
\item $-e_k$: move $-1$ along axis $k$,
\item $0$: stay put.
\end{itemize}
\end{definition}

\begin{definition}[Lattice path]
A lattice path is a finite sequence of steps:
$p = (s_1, s_2, \ldots, s_n)$ with each $s_i$ a lattice step.
\end{definition}

\noindent\textbf{Lean:} \texttt{LatticeStep D} (inductive type),
\texttt{LatticePath D := List (LatticeStep D)}.

\subsection{The winding number}

\begin{definition}[Winding number]
\label{def:winding}
The winding number of path $p$ along axis $k$ is
\[
  w_k(p) := \sum_{i=1}^n \delta_k(s_i),
\]
where $\delta_k(+e_j) = [j = k]$, $\delta_k(-e_j) = -[j = k]$,
$\delta_k(0) = 0$.
\end{definition}

\noindent\textbf{Lean:} \texttt{winding\_number : LatticePath D -> Fin D -> Z}.

\begin{theorem}[Winding number properties {\normalfont(\texttt{WindingCharges})}]
\label{thm:winding-props}
\begin{enumerate}[nosep,label=(\alph*)]
\item $w_k(\emptyset) = 0$ \;(empty path).
\item $w_k([+e_k]) = 1$, \;$w_k([-e_k]) = -1$ \;(unit steps).
\item $w_k([+e_j]) = 0$ for $j \neq k$ \;(orthogonality).
\item $w_k(p_1 \mathbin{+\mkern-6mu+} p_2) = w_k(p_1) + w_k(p_2)$ \;(additivity).
\item If $(s_1, s_2)$ is a cancelling pair ($+e_j$ then $-e_j$), then
  $w_k(p_1 \mathbin{+\mkern-6mu+} [s_1, s_2] \mathbin{+\mkern-6mu+} p_2) = w_k(p_1 \mathbin{+\mkern-6mu+} p_2)$ \;(invariance).
\end{enumerate}
\end{theorem}

\begin{proof}
(a)--(d) by direct computation (\texttt{simp}).
(e) by additivity and the fact that
$\delta_k(+e_j) + \delta_k(-e_j) = 0$.
All proved in Lean, 0~\texttt{sorry}.
\end{proof}

\subsection{Independence}

\begin{theorem}[Winding numbers are independent {\normalfont(\texttt{D\_independent\_charges})}]
\label{thm:independent}
For any $D \geq 1$:
\begin{enumerate}[nosep,label=(\roman*)]
\item There are $D$ independent winding charges.
\item They are pairwise distinguishable: for $j \neq k$, the path
  $[+e_k]$ has $w_j = 0$ and $w_k = 1$.
\item Each charge surjects onto $\Z$: for any $n \in \Z$ and axis $k$,
  there exists a path $p$ with $w_k(p) = n$ and $w_j(p) = 0$ for
  $j \neq k$.
\end{enumerate}
\end{theorem}

\begin{proof}
(i) $|\text{Fin}\,D| = D$.
(ii) Explicit witness.
(iii) Induction on $n$ using \texttt{Int.induction\_on}: append $+e_k$
for the successor case, $-e_k$ for the predecessor case.
All proved in Lean.
\end{proof}

%=============================================================================
\section{Topological Charges}
\label{sec:topo-charge}
%=============================================================================

\begin{definition}[Topological charge]
A topological charge on an $N$-entry ledger is a function
$Q : \text{Configuration}\,N \to \Z$ such that
$Q(\text{next}) = Q(\text{current})$
whenever $\text{next}$ is a variational successor of $\text{current}$.
\end{definition}

\noindent\textbf{Lean:} \texttt{TopologicalCharge N} (structure with
\texttt{value} and \texttt{conserved} fields).

\begin{theorem}[Exact conservation {\normalfont(\texttt{topological\_charge\_trajectory\_conserved})}]
\label{thm:exact-conserved}
For any topological charge $Q$, any variational trajectory
$\text{traj}$, and any tick $t$:
\[
  Q(\text{traj}(t)) = Q(\text{traj}(0)).
\]
\end{theorem}

\begin{proof}
Induction on $t$, using the conservation property at each step.
\end{proof}

\begin{theorem}[Charge algebra {\normalfont(\texttt{addCharges}, \texttt{negCharge})}]
\label{thm:charge-algebra}
Topological charges form an abelian group:
\begin{itemize}[nosep]
\item $(Q_1 + Q_2)(c) := Q_1(c) + Q_2(c)$ is a topological charge.
\item $(-Q)(c) := -Q(c)$ is a topological charge.
\item The zero charge $Q_0(c) := 0$ is a topological charge.
\end{itemize}
\end{theorem}

\begin{theorem}[Noether charges are weaker {\normalfont(\texttt{noether\_not\_necessarily\_quantized})}]
\label{thm:noether-weaker}
There exist Noether charges (conserved $\R$-valued functions) that
are not integer-valued.  Specifically, the log-charge
$\sum_i \log x_i$ takes the non-integer value $\log 2$ on the
configuration $x_i = 2$ for all~$i$.

\emph{Proof:} $0 < \log 2 < 1$, so $\log 2$ is not an integer.
(The proof that $\log 2 < 1$ uses $2 < e$ and the monotonicity of $\log$.)
\end{theorem}

%=============================================================================
\section{Linking Numbers in $D = 3$}
\label{sec:linking}
%=============================================================================

\subsection{The Hopf link on the cubic lattice}

\begin{definition}[Closed path]
A lattice path is closed if $w_k(p) = 0$ for all $k$.
\end{definition}

We construct an explicit linked pair of closed paths in $\Z^3$:

\begin{definition}[Hopf link on $\Z^3$]
\label{def:hopf}
\begin{align*}
C_1 &:= [+e_x,\; +e_y,\; -e_x,\; -e_y]
  &&\text{(square in the $xy$-plane)}\\
C_2 &:= [+e_z,\; +e_x,\; -e_z,\; -e_x]
  &&\text{(square in the $xz$-plane, shifted)}
\end{align*}
\end{definition}

\begin{theorem}[$C_1$ and $C_2$ are closed {\normalfont(\texttt{hopf\_1\_closed}, \texttt{hopf\_2\_closed})}]
Both paths have zero winding number on all 3 axes.
\end{theorem}

\begin{proof}
Each axis gets one $+$ step and one $-$ step.
Verified by \texttt{fin\_cases k <;> simp} in Lean.
\end{proof}

\begin{theorem}[The Hopf link is linked {\normalfont(\texttt{hopf\_link\_exists\_D3})}]
\label{thm:hopf-linked}
$C_2$ passes through the interior of the square bounded by $C_1$.
Specifically, $C_2$'s $+e_z$ step crosses the $z = 0$ plane at position
$(0, 0, 0)$, which is inside the region $\{0 \leq x \leq 1,\,
0 \leq y \leq 1\}$ bounded by $C_1$.
\end{theorem}

\subsection{No linking in low dimensions}

\begin{theorem}[No linking in $D = 1$ {\normalfont(\texttt{D1\_no\_linking\_structure})}]
Every closed path in $\Z^1$ has zero winding number.
Two closed paths in $\Z^1$ cannot link because $\Z$ is totally ordered:
one path is always ``to the left'' of the other.
\end{theorem}

\begin{theorem}[No linking in $D = 2$ {\normalfont(Jordan curve theorem, informal)}]
In $\Z^2$, any closed non-self-intersecting path bounds a simply connected
region.  A second closed path is either entirely inside or entirely outside
this region.  There is no way to link.
\end{theorem}

\subsection{The dimension dependence}

\begin{theorem}[Independent loops {\normalfont(\texttt{independent\_loop\_count})}]
\label{thm:loop-count}
In $D$ dimensions, the number of independent non-trivial square loops is
$\binom{D}{2} = D(D-1)/2$.  For $D = 3$: $\binom{3}{2} = 3$.
\end{theorem}

\begin{center}
\begin{tabular}{@{}cccl@{}}
\toprule
$D$ & Indep.\ loops & Linking? & Conservation charges \\
\midrule
1 & 0 & No & 0 \\
2 & 1 & No (Jordan) & 0 \\
\textbf{3} & \textbf{3} & \textbf{Yes (Hopf)} & \textbf{3} \\
4 & 6 & Trivial (unlinks) & 0 \\
\bottomrule
\end{tabular}
\end{center}

%=============================================================================
\section{The Three Charges of the Standard Model}
\label{sec:sm-charges}
%=============================================================================

\subsection{The identification}

\begin{theorem}[All the ``3''s are the same {\normalfont(\texttt{all\_threes\_unified})}]
\label{thm:all-threes}
The following quantities are all equal to~$3$:
\begin{enumerate}[nosep]
\item Number of independent winding charges in $\Z^3$.
\item Number of face-pairs of the cube $Q_3$.
\item Number of quark colors $N_c$.
\item Number of SM conserved charges (electric, baryon, lepton).
\item Number of fermion generations.
\end{enumerate}
\end{theorem}

\begin{proof}
All reduce to $D = 3$ by definition or by the dimension-forcing theorem.
Proved in Lean by \texttt{rfl} and \texttt{Fintype.card\_fin 3}.
\end{proof}

\subsection{The physical assignment}

\begin{definition}[Charge-to-axis bijection {\normalfont(\texttt{charge\_to\_axis})}]
\begin{align*}
\text{Electric charge} &\;\leftrightarrow\; e_x \text{ (axis 0)} \\
\text{Baryon number} &\;\leftrightarrow\; e_y \text{ (axis 1)} \\
\text{Lepton number} &\;\leftrightarrow\; e_z \text{ (axis 2)}
\end{align*}
\end{definition}

\begin{theorem}[Bijection {\normalfont(\texttt{charge\_to\_axis\_bijective})}]
The map $\text{SMCharge} \to \text{Fin}\,3$ is a bijection
(injective + surjective, proved by case analysis in Lean).
\end{theorem}

\subsection{Why topological charges are quantized}

Charge quantization is the observation that all observed electric
charges are integer multiples of $e/3$.  In the standard framework,
this requires the Dirac quantization argument (magnetic monopoles)
or grand unification.

In RS, quantization is \emph{automatic}: winding numbers are integers
by construction.  The codomain is $\Z$, not $\R$.  There is nothing
to explain---the charge cannot take non-integer values because it counts
lattice steps.

\subsection{Why topological charges are absolutely conserved}

Baryon and lepton number conservation are ``accidental'' symmetries
in the Standard Model---they are not gauged and can in principle be
violated by non-perturbative effects (sphalerons, proton decay in GUTs).

In RS, these charges are \emph{topological}, not accidental.  A linking
number cannot change under continuous deformation.  On the lattice,
``continuous deformation'' means inserting/removing cancelling pairs,
which we have proved preserves all winding numbers.  Therefore:

\begin{theorem}[Absolute conservation]
Topological charges cannot be violated by any process in the variational
dynamics.  There are no anomalies, no sphalerons, no proton decay.
\end{theorem}

\begin{remark}[Falsifiable prediction]
The absolute conservation of baryon number is a falsifiable prediction:
if proton decay is observed, RS is falsified.  Conversely, the
continued non-observation of proton decay (current bound:
$\tau_p > 10^{34}$ years) is consistent with topological conservation
and inconsistent with many GUT predictions.
\end{remark}

%=============================================================================
\section{The Universe Starts Neutral}
\label{sec:neutral}
%=============================================================================

\begin{theorem}[Initial neutrality {\normalfont(\texttt{total\_charge\_always\_zero})}]
If the universe begins with zero topological charge (as forced by the
Past Theorem: the initial state is the unity configuration with all
entries $= 1$), then the total charge remains zero at all future ticks.
\end{theorem}

\begin{proof}
$Q(\text{traj}(t)) = Q(\text{traj}(0)) = 0$ for all $t$, by
Theorem~\ref{thm:exact-conserved}.
\end{proof}

\begin{corollary}
The universe has zero total electric charge, zero total baryon number,
and zero total lepton number.  Any positive charge is balanced by an
equal negative charge elsewhere in the ledger.
\end{corollary}

%=============================================================================
\section{Conclusion}
\label{sec:conclusion}
%=============================================================================

We have derived the three conservation laws of the Standard Model from
the topology of the cubic lattice $\Z^3$, without assuming any
continuous symmetry.

The key insight is that conservation laws are not imposed by Noether's
theorem (continuous symmetry $\to$ conserved current).  They are imposed
by \emph{topology}: the winding numbers of lattice paths are integer-valued
invariants that cannot be changed by local deformations.

The derivation chain is:
\[
\boxed{
\text{RCL} \to \Jcost \text{ unique} \to D = 3 \text{ forced}
\to \Z^3 \text{ lattice} \to 3 \text{ winding numbers}
\to 3 \text{ conserved charges}
}
\]

Every link is proved in Lean~4.  The physical identification
(electric charge, baryon number, lepton number) is a bijection between
$\text{SMCharge}$ and $\text{Fin}\,3$, proved injective and surjective
by case analysis.

This represents a fundamentally different origin for conservation laws
than anything in the standard framework.  It explains charge
quantization (winding numbers are integers), it predicts absolute
conservation (no proton decay), and it unifies the ``3'' in the number
of conserved charges with the ``3'' in the spatial dimension, the number
of quark colors, and the number of fermion generations.

All of these are the same ``3'': $\dim \Z^3 = 3$.

\bigskip
\begin{center}
\textsc{--- End ---}
\end{center}

\end{document}
